\documentclass[10pt,a4paper]{amsart}
\usepackage[margin=19mm]{geometry}
\usepackage{amsmath,amssymb,mathtools}
\usepackage[scr=rsfs,cal=euler]{mathalfa}
\usepackage{xfrac}
\usepackage[unicode,hidelinks,bookmarks=false]{hyperref}
\newcommand{\ie}{i.e.\ }
\newcommand{\cf}{cf.\ }
\newcommand{\resp}{respectively\ }
\newcommand{\Subsection}{\subsection}
\providecommand{\nobreakdash}{\nobreak\hbox{-}}
\setlength{\emergencystretch}{2em}
\multlinegap0pt
\usepackage{amssymb}
\newtheorem{lemma}{Lemma}
\newtheorem{proposition}{Proposition}
\title{Homological framework of noncommutative complex analytic geometry and functional calculus}
\author{Anar Dosi}
\date{Source: arXiv:2512.20191v1 (CC BY 4.0)}
\begin{document}
\maketitle

\noindent\emph{Source and licence.} Homological framework of noncommutative complex analytic geometry and functional calculus. Version arXiv:2512.20191v1 (CC BY 4.0), distributed by arXiv under the Creative Commons Attribution 4.0 International licence, http://creativecommons.org/licenses/by/4.0/, as recorded in the arXiv licence metadata for this exact version and shown on the abstract page https://arxiv.org/abs/2512.20191v1. Full licence notice: \texttt{licenses/arxiv-2512.20191-CC-BY-4.0.txt}.\par\smallskip

\noindent\textit{Editorial scope.} Counted unit: the article's proposition corCes1 (source0.tex lines 1364--1367). The article's complete original proof of it is delivered with it (source0.tex lines 1369--1442, one \texttt{proof} environment of 74 lines). No mathematics is written, repaired or re-derived: the statement and the proof are the article's own lines, in the article's own notation, and no retained line is substituted or re-flowed.  Retained uncounted as the source-local context the proof needs: lines 1228--1233; lines 1310--1315.  Nothing was omitted from any retained span, so the package carries no omission record.  No retained line contains a citation, so the wrapper carries no bibliography and no bibliography entry.  No retained line contains a figure environment, an image-inclusion command or drawing code, so no image is delivered and the package declares no asset dependency.  Editorial formatting only, in addition: an AMS \texttt{amsart} preamble that restates the article's own declarations and macros used by the retained lines, namely the environments \texttt{lemma}, \texttt{proposition} on separate counters, together with the standard packages those lines need.  The retained environments print in this document as Lemma 1, Proposition 1 in order of appearance, whereas the article prints the same items with the numbering of its own full document.  The complete native source of the article as distributed in the arXiv e-print for this exact version is bundled unchanged as \texttt{sources/191545/source0.tex}.
\begin{lemma}
\label{lemCes1}A poset $A$-category $\mathcal{S}$ is a complete lattice if and
only if there exists infimum $\wedge U$ for every open subset $U\subseteq
\mathcal{S}$ bounded below. In this case, $\mathcal{S}$ is an irreducible
topological space.
\end{lemma}

\begin{lemma}
\label{lemCes2}If $F:\mathcal{S}\rightarrow\mathcal{G}$ is a morphism of the
category $\mathfrak{C}_{\Omega}$, then it defines a morphism of the presheaves
$\mathcal{S}^{\tau}\rightarrow\mathcal{G}^{\tau}$ on $\Omega$. Thus we have
the functor $\mathfrak{C}_{\Omega}\rightarrow\mathfrak{Fa}_{\Omega}$.
\end{lemma}

\begin{proposition}
\label{corCes1}The functor $\mathfrak{C}_{\Omega}\rightarrow\mathfrak{Fa}%
_{\Omega}$ is a category equivalence.
\end{proposition}

\begin{proof}
By Lemma \ref{lemCes2}, we have the well defined functor $\mathfrak{C}%
_{\Omega}\rightarrow\mathfrak{Fa}_{\Omega}$, $\mathcal{S\mapsto S}^{\tau}$.
Let $\mathcal{P}:\Omega^{\operatorname{op}}\rightarrow\mathfrak{Fa}$ be a
Fr\'{e}chet algebra presheaf on $\Omega$ with the Fr\'{e}chet algebra $A$ of
the global sections. Put $\widehat{\mathcal{P}}$ to be a poset category
defined by the presheaf $\mathcal{P}$, that is, the objects of
$\widehat{\mathcal{P}}$ are the algebras $\mathcal{P}\left(  V\right)  $ over
open subsets $V\in\Omega$ with their unique morphisms $\mathcal{P}\left(
V\right)  \rightarrow\mathcal{P}\left(  W\right)  $ given by the restriction
morphisms of the presheaf $\mathcal{P}$. One can easily verify that
$\mathcal{P}\left(  V\right)  \vee\mathcal{P}\left(  W\right)  =\mathcal{P}%
\left(  V\cap W\right)  $ and $\mathcal{P}\left(  V\right)  \wedge
\mathcal{P}\left(  W\right)  =\mathcal{P}\left(  V\cup W\right)  $ \ Moreover,
every subset $\left\{  \mathcal{P}\left(  V_{i}\right)  \right\}  $ of
$\widehat{\mathcal{P}}$ is bounded below by $A$ and $\wedge\left\{
\mathcal{P}\left(  V_{i}\right)  \right\}  =\mathcal{P}\left(  \cup_{i}%
V_{i}\right)  $. Hence $\widehat{\mathcal{P}}$ is a unital complete lattice
$A$-category. An open subset $U^{\prime}\subseteq\widehat{\mathcal{P}}$
corresponds to some $U\in\Omega$, that is, it consists of those $\mathcal{P}%
\left(  V\right)  $ with $V\subseteq U$. Thus $\widehat{\mathcal{P}}_{\tau}$
consists of $U^{\prime}=\left\{  \mathcal{P}\left(  V\right)  :V\subseteq
U\right\}  $ with $U\in\Omega$. In this case, we have $\wedge U^{\prime
}=\mathcal{P}\left(  U\right)  $ (see Lemma \ref{lemCes1}), and the
correspondence $\mathcal{P}\longrightarrow\widehat{\mathcal{P}}$ is
functorial. Namely, assume that $\mathcal{P}_{1}\rightarrow\mathcal{P}_{2}$ is
a presheaf morphism over $\Omega$, which is a functor transformation
\[
\left(  \mathcal{P}_{1}:\Omega^{\operatorname{op}}\rightarrow\mathfrak{Fa}%
\right)  \rightarrow\left(  \mathcal{P}_{2}:\Omega^{\operatorname{op}%
}\rightarrow\mathfrak{Fa}\right)  .
\]
If $\widehat{\mathcal{P}_{i}}$, $i=1,2$ are the related unital complete
lattice categories, then%
\[
F:\widehat{\mathcal{P}_{1}}\rightarrow\widehat{\mathcal{P}_{2}},\quad
\mathcal{P}_{1}\left(  V\right)  \rightarrow F\left(  \mathcal{P}_{1}\left(
V\right)  \right)  =\mathcal{P}_{2}\left(  V\right)  ,\quad V\in\Omega
\]
is a morphism of these categories. Actually, $F$ is a strong morphism. Indeed,
if $U\in\Omega$ with its open subset representations $U_{i}=\left\{
\mathcal{P}_{i}\left(  V\right)  :V\subseteq U\right\}  $ in
$\widehat{\mathcal{P}_{i}}$, then
\[
F\left(  U_{1}\right)  =F\left\{  \mathcal{P}_{1}\left(  V\right)  :V\subseteq
U\right\}  =\left\{  \mathcal{P}_{2}\left(  V\right)  :V\subseteq U\right\}
=U_{2}.
\]
Thus $F:\left(  \widehat{\mathcal{P}_{1}}\right)  _{\tau}\rightarrow\left(
\widehat{\mathcal{P}_{2}}\right)  _{\tau}$ is a category equivalence of the
related poset categories of open subsets, which means that $F$ is a strong
morphism. Hence we come up with a functor $\mathfrak{Fa}_{\Omega}%
\rightarrow\mathfrak{C}_{\Omega}$.

If $\mathcal{S}$ is a unital complete lattice $A$-category with its
corresponding presheaf $\mathcal{S}^{\tau}$ on $\Omega$, then $\mathcal{S}%
^{\tau}\left(  U_{\mathcal{A}}\right)  =\wedge U_{\mathcal{A}}=\mathcal{A}$
for every object $\mathcal{A}$ of $\mathcal{S}$. Moreover, $\mathcal{A\leq B}$
in $\mathcal{S}$ iff $U_{\mathcal{A}}\supseteq U_{\mathcal{B}}$, which means
that there is a morphism $\mathcal{S}^{\tau}\left(  U_{\mathcal{A}}\right)
\rightarrow\mathcal{S}^{\tau}\left(  U_{\mathcal{B}}\right)  $. Thus the poset
category $\widehat{\mathcal{S}^{\tau}}$ of the presheaf $\mathcal{S}^{\tau}$
is reduced to $\mathcal{S}$ up to an order isomorphism. Conversely, if
$\mathcal{P}$ is a presheaf over $\Omega$ with its corresponding poset
category $\widehat{\mathcal{P}}$, then $\widehat{\mathcal{P}}$ is a unital
complete lattice $A$-category. Let $\left(  \widehat{\mathcal{P}}\right)
^{\tau}$ be the related presheaf on the irreducible topological space
$\widehat{\mathcal{P}}$ (see Lemma \ref{lemCes1}). Then the category of the
open subsets of $\widehat{\mathcal{P}}$ is reduced to $\Omega$, and $\left(
\widehat{\mathcal{P}}\right)  ^{\tau}\left(  U^{\prime}\right)  =\wedge
U^{\prime}=\mathcal{P}\left(  U\right)  $ for an open subset $U^{\prime
}\subseteq\widehat{\mathcal{P}}$ corresponding to $U\in\Omega$. Thus the
result follows.
\end{proof}

\end{document}
